% !TeX root = ../../tesis.tex
%%%%%%%%%%%%%%%%%%%%%%%%%%%%%%%%%%%%%%%%%%%%%%%%%%%%%%%%%%%%%%%%%%%
% SECCIÓN 5.5.1 — TIEMPO DE VIAJE PROMEDIO (T)
%%%%%%%%%%%%%%%%%%%%%%%%%%%%%%%%%%%%%%%%%%%%%%%%%%%%%%%%%%%%%%%%%%%
\subsection{Tiempo de Viaje Promedio (\texorpdfstring{$T$}{T})}
\label{subsec:tiempo-viaje-promedio}

Ejecutado el algoritmo de \dijkstra{} ponderado sobre el componente
gigante, el tiempo de viaje promedio de la red resulta en
$T = 108.92$ minutos (corrida: 2026-09-02, 10,561~nodos, tiempo de
cómputo: 251.8~segundos). El valor sintetiza el costo temporal medio
de todos los pares origen-destino posibles bajo las condiciones de
ponderación de Fase~2.

La referencia de campo más cercana proviene de la Encuesta Nacional de
Ocupación y Empleo (ENOE), que estima en aproximadamente 85~minutos el
tiempo promedio de traslado en la \zmvm{}. Frente a los 108.92~minutos
del modelo, la brecha de $\sim$24~minutos tiene origen metodológico, no
en una sobreestimación del sistema. Mientras la ENOE construye su
promedio a partir de desplazamientos observados —concentrados en los
corredores de alta afluencia que conectan el centro con la periferia
próxima— el modelo asigna el mismo peso a cada par del componente
gigante. Los trayectos entre periferias opuestas, que la demanda
cotidiana apenas realiza, quedan contabilizados en igualdad de
condiciones que los recorridos más frecuentes, elevando el promedio
topológico por encima del promedio ponderado por demanda.

Esa variación queda desglosada por alcaldía en la
Figura~\ref{fig:grafica-t-promedio}.

\begin{figure}[H]
    \centering
    \includegraphics[width=0.92\textwidth]{%
        Figures/Mapas/Tableau_Resources/RedActual/VFTModel/UniGrafica/%
        Grafica_VFTModel_2026_Tiempo_Promedio}
    \caption[Tiempo de viaje promedio por alcaldía --- \zmvm{}]{%
    Tiempo de viaje promedio ($T$) por alcaldía. Corrida sobre el
    componente gigante (10,561~nodos) con \dijkstra{} ponderado.
    \zmvm{}, 2026-09-02.}
    \label{fig:grafica-t-promedio}
    \fuentefigura{Elaboración propia con \textbf{Tableau};
    datos: \texttt{VFTModel} \citep{cabrera2026VFTModel}.}
\end{figure}

Con $T$ calculado, surge una pregunta estructural: ¿qué nodos
concentran mayor responsabilidad sobre ese valor? No todos los nodos
son equivalentes en la topología del grafo; algunos articulan una
fracción desproporcionada de las rutas mínimas y su fallo impacta los
tiempos en toda la red. La centralidad de intermediación $B(v)$
—analizada en la sección siguiente— identifica esos puntos.
